\documentclass[10pt,a4paper]{amsart}
\usepackage{amsmath,amssymb,amsthm,graphicx}
% Retained minimal preamble for the Stacks Project chapter more-morphisms.tex (More on Morphisms).
% Verbatim excerpt of the chapter preamble: Xy-pic setup (native preamble lines 14-18) and the native
% theorem declarations with the mathematical macros (native preamble lines 154-217). Copied mechanically;
% nothing added, removed or redefined. The wrapper class and the amsmath/amssymb/amsthm/graphicx base
% packages are supplied by the assembler, so the class/externaldocument part of the native preamble is omitted.
% For commutative diagrams we use Xy-pic
\usepackage[all]{xy}

% We use 2cell for 2-commutative diagrams.
\xyoption{2cell}
% Theorem environments.
%
\theoremstyle{plain}
\newtheorem{theorem}[subsection]{Theorem}
\newtheorem{proposition}[subsection]{Proposition}
\newtheorem{lemma}[subsection]{Lemma}

\theoremstyle{definition}
\newtheorem{definition}[subsection]{Definition}
\newtheorem{example}[subsection]{Example}
\newtheorem{exercise}[subsection]{Exercise}
\newtheorem{situation}[subsection]{Situation}

\theoremstyle{remark}
\newtheorem{remark}[subsection]{Remark}
\newtheorem{remarks}[subsection]{Remarks}

\numberwithin{equation}{subsection}

% Macros
%
\def\lim{\mathop{\mathrm{lim}}\nolimits}
\def\colim{\mathop{\mathrm{colim}}\nolimits}
\def\Spec{\mathop{\mathrm{Spec}}}
\def\Hom{\mathop{\mathrm{Hom}}\nolimits}
\def\Ext{\mathop{\mathrm{Ext}}\nolimits}
\def\SheafHom{\mathop{\mathcal{H}\!\mathit{om}}\nolimits}
\def\SheafExt{\mathop{\mathcal{E}\!\mathit{xt}}\nolimits}
\def\Sch{\mathit{Sch}}
\def\Mor{\mathop{\mathrm{Mor}}\nolimits}
\def\Ob{\mathop{\mathrm{Ob}}\nolimits}
\def\Sh{\mathop{\mathit{Sh}}\nolimits}
\def\NL{\mathop{N\!L}\nolimits}
\def\CH{\mathop{\mathrm{CH}}\nolimits}
\def\proetale{{pro\text{-}\acute{e}tale}}
\def\etale{{\acute{e}tale}}
\def\QCoh{\mathit{QCoh}}
\def\Ker{\mathop{\mathrm{Ker}}}
\def\Im{\mathop{\mathrm{Im}}}
\def\Coker{\mathop{\mathrm{Coker}}}
\def\Coim{\mathop{\mathrm{Coim}}}

% Boxtimes
%
\DeclareMathSymbol{\boxtimes}{\mathbin}{AMSa}{"02}

%
% Macros for moduli stacks/spaces
%
\def\QCohstack{\mathcal{QC}\!\mathit{oh}}
\def\Cohstack{\mathcal{C}\!\mathit{oh}}
\def\Spacesstack{\mathcal{S}\!\mathit{paces}}
\def\Quotfunctor{\mathrm{Quot}}
\def\Hilbfunctor{\mathrm{Hilb}}
\def\Curvesstack{\mathcal{C}\!\mathit{urves}}
\def\Polarizedstack{\mathcal{P}\!\mathit{olarized}}
\def\Complexesstack{\mathcal{C}\!\mathit{omplexes}}
% \Pic is the operator that assigns to X its picard group, usage \Pic(X)
% \Picardstack_{X/B} denotes the Picard stack of X over B
% \Picardfunctor_{X/B} denotes the Picard functor of X over B
\def\Pic{\mathop{\mathrm{Pic}}\nolimits}
\def\Picardstack{\mathcal{P}\!\mathit{ic}}
\def\Picardfunctor{\mathrm{Pic}}
\def\Deformationcategory{\mathcal{D}\!\mathit{ef}}
\usepackage{cleveref}
\crefname{originaldefinition}{original definition}{original definitions}
\Crefname{originaldefinition}{Original definition}{Original definitions}
\makeatletter
\newlabel{definition-thickening}{{04EX}{source}{}{sourcepointer-1}{}}
\newlabel{definition-thickening@cref}{{[originaldefinition][0][]04EX}{[1][0][]source}{}{}{}}
\makeatother
\crefname{originallemma}{original lemma}{original lemmas}
\Crefname{originallemma}{Original lemma}{Original lemmas}
\makeatletter
\newlabel{lemma-picard-group-first-order-thickening}{{0C6R}{source}{}{sourcepointer-2}{}}
\newlabel{lemma-picard-group-first-order-thickening@cref}{{[originallemma][0][]0C6R}{[1][0][]source}{}{}{}}
\makeatother
\crefname{originallemma}{original lemma}{original lemmas}
\Crefname{originallemma}{Original lemma}{Original lemmas}
\makeatletter
\newlabel{algebra-lemma-lift-nth-roots}{{0CAP}{source}{}{sourcepointer-3}{}}
\newlabel{algebra-lemma-lift-nth-roots@cref}{{[originallemma][0][]0CAP}{[1][0][]source}{}{}{}}
\makeatother
\title{The n-torsion in the Picard group is insensitive to a thickening when n is invertible}
\author{The Stacks Project Authors}
\date{stacks-project:more-morphisms.tex:local-snapshot-20261011; GNU FDL 1.2 or later}
\begin{document}
\maketitle
\noindent\textit{Editorial note.} All delivered mathematical text is the author's own, copied from the pinned Stacks Project chapter file more-morphisms.tex; no statement or proof step was rewritten, reordered or removed. The author gives two proof environments for this one lemma, headed 'Proof for a finite order thickening' and 'Proof in general'; both are included in the counted spans, in their original order, so the delivered material carries the complete author proof of the statement.\par\medskip
\noindent\textit{Editorial note.} The preamble is the chapter's own retained minimal preamble (sources/stacks-more-morphisms/retained-preamble.tex), a verbatim excerpt of the native Xy-pic setup and of the native theorem declarations and mathematical macros; no macro was invented or redefined.\par\medskip
\noindent\textit{Editorial note.} The two same-chapter references are located at their \textbackslash{}label lines inside the pinned chapter file; the cross-chapter pointer is located in the project's own tag list staged in the snapshot (sources/stacks-chow/tags.txt), at the verbatim line carrying its canonical tag.\par\medskip
\noindent\textit{Allowed standard prerequisites (editorial).}\begin{itemize}
\item Thickening, and the reduction from finite order thickenings to first order thickenings. Same chapter, Definition environment labelled definition-thickening (source line 36) followed by the general principle stated immediately after it; invoked at proof lines 640-641 and again in the general proof at line 652. Supplies the notion of thickening and the reduction principle that lets the author treat only first order thickenings in the first proof, and gives the ideal sheaf I cutting out X in X' in the second.
\item The long exact sequence for the Picard group of a first order thickening. Same chapter, Lemma labelled lemma-picard-group-first-order-thickening (source line 599, Stacks tag 0C6R, counted separately in this night's work); invoked at proof line 643 and again at line 659. Supplies the exact sequence relating Pic(X'), Pic(X) and the cohomology of the ideal sheaf, which is the spine of both proofs of this lemma.
\item Lifting nth roots of units congruent to one modulo a square-zero ideal. Algebra, Lemma labelled algebra-lemma-lift-nth-roots (canonical Stacks tag 0CAP); invoked at proof line 664. Supplies that raising to the nth power is an isomorphism on the group of sections one plus a locally nilpotent sheaf of ideals; the assembly of the two proofs is retained in the counted spans.
\end{itemize}\medskip
\ifdefined\hypertarget\hypertarget{sourcepointer-1}{}\fi
\noindent\textit{Source reference.} Original definition 04EX: sources/stacks-more-morphisms/pinned-original-source.tex, source lines 36--36. Same-chapter pointer invoked at proof line 641 for the reduction principle and referred to again in the general proof at line 652; both proofs are retained verbatim in the counted spans.\par\medskip
\ifdefined\hypertarget\hypertarget{sourcepointer-2}{}\fi
\noindent\textit{Source reference.} Original lemma 0C6R: sources/stacks-more-morphisms/pinned-original-source.tex, source lines 599--599. Same-chapter pointer invoked at proof lines 643 and 659 for the exact sequence relating Pic(X') and Pic(X) to the cohomology of the ideal sheaf; both uses are inputs whose consequences are drawn inside the counted spans.\par\medskip
\ifdefined\hypertarget\hypertarget{sourcepointer-3}{}\fi
\noindent\textit{Source reference.} Original lemma 0CAP: sources/stacks-chow/tags.txt, source lines 15444--15444. Cross-chapter standard prerequisite invoked once at proof line 664 for invertibility of the nth power map on units congruent to one; the reduction of the general case to that statement is retained in the counted span.\par\medskip
\begin{lemma}
\label{lemma-torsion-pic-thickening}
Let $X \subset X'$ be a thickening. Let $n$ be an integer
invertible in $\mathcal{O}_X$. Then the map
$\Pic(X')[n] \to \Pic(X)[n]$ is bijective.
\end{lemma}

\begin{proof}[Proof for a finite order thickening]
By the general principle explained following
Definition \ref{definition-thickening}
this reduces to the case of a first order thickening.
Then may use Lemma \ref{lemma-picard-group-first-order-thickening}
to see that it suffices to show that
$H^1(X, \mathcal{I})[n]$, $H^1(X, \mathcal{I})/n$, and
$H^2(X, \mathcal{I})[n]$ are zero.
This follows as multiplication by $n$ on $\mathcal{I}$
is an isomorphism as it is an $\mathcal{O}_X$-module.
\end{proof}

\begin{proof}[Proof in general]
Let $\mathcal{I} \subset \mathcal{O}_{X'}$ be the quasi-coherent ideal
sheaf cutting out $X$. Then we have a short exact sequence of
abelian groups
$$
0 \to (1 + \mathcal{I})^* \to \mathcal{O}_{X'}^* \to \mathcal{O}_X^* \to 0
$$
We obtain a long exact cohomology sequence as in the statement of
Lemma \ref{lemma-picard-group-first-order-thickening}
with $H^i(X, \mathcal{I})$ replaced by $H^i(X, (1 + \mathcal{I})^*)$.
Thus it suffices to show that raising to the $n$th power is an
isomorphism $(1 + \mathcal{I})^* \to (1 + \mathcal{I})^*$.
Taking sections over affine opens this follows from
Algebra, Lemma \ref{algebra-lemma-lift-nth-roots}.
\end{proof}

\end{document}
